\chapter{From Energy-Based Models to Diffusion and Flows}
\label{ch:diffusion}\label{ch:ebm}

Chapter~\ref{ch:statmech} left the generative branch of machine learning
at the Boltzmann machine: a model is an energy, learning is moment
matching, sampling is dynamics. This chapter follows that branch to the
present. It reviews, in the depth their role in this thesis demands, the
three ideas the research chapters build on: the energy-based view of
modelling \citep{lecun2006tutorial}, the score-matching principle that
removed the partition function from training
\citep{hyvarinen2005estimation, vincent2011connection}, and the diffusion
construction that turned denoising into a complete generative method
\citep{sohl2015deep, ho2020denoising, song2021scorebased}, together with
its deterministic sibling, the flow model. The chapter then reviews the
statistical-physics literature on diffusion models, from the
high-dimensional analyses of \citet{biroli2023generative} to the
speciation, memorisation, and generalisation results of the
Biroli--M\'ezard line of work, and, on the training side, the cascade of
phase transitions in energy-based-model learning
\citep{bachtis2024cascade}. It closes by stating the gap this thesis
addresses.

\section{Energy-based models}\label{sec:ebm-def}

Modern generative modelling asks a single question: given samples of an
unknown distribution $p_0$, build a machine that produces new ones. The
last decade produced several answers: variational autoencoders
\citep{kingma2014auto}, adversarial networks
\citep{goodfellow2014generative}, normalizing flows
\citep{rezende2015variational}, diffusion models. It is easy to present
them as unrelated tricks. The energy-based view of
\citet{lecun2006tutorial} is the antidote: it is the direct continuation
of the Boltzmann-machine programme of
Section~\ref{sec:boltzmann-machines}, and, as this chapter unfolds,
diffusion models will appear as the member of the energy-based family
that finally solved its oldest problem.

An \emph{energy-based model} (EBM) over $x \in \mathcal{X}$ assigns to
each configuration a scalar energy $E_\theta(x)$ measuring how badly $x$
fits the model's idea of the data: low energy for plausible
configurations, high for implausible ones. In
\citeauthor{lecun2006tutorial}'s formulation the energy is the primitive
object: inference means finding low-energy configurations, and learning
means shaping the landscape, pushing the energy of observed data down and
the energy of everything else up. A probabilistic model is recovered
whenever one normalises with the Gibbs measure of
Chapter~\ref{ch:statmech},
\begin{equation}
  p_\theta(x) \;=\; \frac{e^{-E_\theta(x)}}{Z(\theta)},
  \qquad
  Z(\theta) = \int_{\mathcal{X}} e^{-E_\theta(x)}\, \dd x ,
  \label{eq:ebm}
\end{equation}
the Boltzmann distribution \eqref{eq:boltzmann} with $\beta$ absorbed
into the energy and the physics replaced by a parametric family. The
appeal of \eqref{eq:ebm} is that $E_\theta$ is unconstrained: any
function defines a valid unnormalised density. The price is the
partition function, and every basic operation hits it. \emph{Evaluating}
$p_\theta(x)$ requires $Z(\theta)$, intractable in general
(Section~\ref{sec:ising}). \emph{Sampling} requires the MCMC machinery of
Section~\ref{sec:mcmc}, whose mixing time can be exponential in the
presence of metastability. \emph{Learning} by maximum likelihood requires
the gradient
\begin{equation}
  -\nabla_\theta \log p_\theta(x)
  = \nabla_\theta E_\theta(x)
  - \E_{p_\theta}\!\big[ \nabla_\theta E_\theta \big],
  \label{eq:ebm-gradient}
\end{equation}
the ``push down on data, pull up on the model's own samples'' rule, the
continuous-variable form of the Boltzmann-machine rule
\eqref{eq:bm-learning}, whose second term again demands sampling from
$p_\theta$. \citet{lecun2006tutorial} frame the alternatives
systematically: either fight the negative phase with \emph{contrastive}
losses that pull up only on well-chosen offending configurations
(contrastive divergence \citep{hinton2002training} being the
MCMC-flavoured member of the family), or design architectures and
regularisers that bound the volume of low-energy space so that no pull-up
is needed. The lesson of the tutorial is that the partition function is
not an accounting detail; it is \emph{the} design constraint of the
field.

There are two classical escape routes, and this thesis walks both. The
first is the \emph{score}: since
$\nabla_x \log p_\theta(x) = -\nabla_x E_\theta(x)$, the gradient of the
log-density in $x$ is free of $Z$, because differentiating in $x$ kills
the constant. Sections~\ref{sec:score-matching}--\ref{sec:tweedie} build
the modern edifice standing on this observation. The second is
\emph{structure}: when the energy decomposes into local terms, marginals
and even $Z$ itself become computable by dynamic programming; that route
(Markov random fields, factor graphs, message passing) is developed in
Chapter~\ref{ch:graphs}.

\section{The diffusion idea}\label{sec:diffusion-idea}

Diffusion models \citep{sohl2015deep, ho2020denoising,
song2021scorebased} solve the generative problem with a two-process
design inspired by non-equilibrium thermodynamics: a fixed \emph{forward
process} gradually corrupts data into pure noise, and a learned
\emph{reverse process} undoes the corruption step by step. A recent
book-length survey of the family is \citet{lai2025principles}. The
forward process used throughout this thesis is the OU channel of
Section~\ref{sec:ou}: at diffusion time $t$,
\begin{equation}
  x_t \;=\; a\, e^{-t} + \sqrt{\Deltat}\;\varepsilon,
  \qquad a \sim p_0,\quad \varepsilon \sim \Nd(0, I),
  \qquad \Deltat = 1 - e^{-2t},
  \label{eq:forward-marginal}
\end{equation}
so that the noised marginal is the convolution
\begin{equation}
  p_t(x) \;=\; \int p_0(a)\; \Nd\!\big(x;\, a e^{-t},\, \Deltat I\big)\,
  \dd a .
  \label{eq:pt-convolution}
\end{equation}
As $t$ grows, $p_t$ interpolates from the data distribution to
$\Nd(0, I)$. By Anderson's theorem (Section~\ref{sec:reversal}), the
corruption is undone by the reverse SDE \eqref{eq:reverse-sde}, whose
only unknown ingredient is the score $\nabla_x \log p_t(x)$ of the noised
marginals. Generative diffusion is therefore, at bottom, \emph{score
estimation}: whoever knows $\nabla \log p_t$ for all $t$ can sample from
$p_0$. Note what this design buys relative to
Section~\ref{sec:ebm-def}: no energy is ever normalised, no equilibrium
is ever awaited. The Markov-chain convergence problem of MCMC is replaced
by a \emph{fixed-horizon} integration of an SDE whose drift is learned.
The partition-function problem has not been solved; it has been
side-stepped, by working only with objects (scores) from which $Z$ has
already been differentiated away.

\subsection{The thermodynamic reading, and the direction of the arrow}
\label{sec:diffusion-thermo}

The appeal to non-equilibrium thermodynamics in \citet{sohl2015deep} is
structural rather than decorative. Their forward process is a chain of small
perturbations converting a structured distribution into a simple one, that is
an annealing schedule, and the objective they derive is a variational bound of
the same algebraic form as the work identities of non-equilibrium statistical
mechanics: the log-likelihood is bounded by an average of log-ratios along the
forward path, exactly as free-energy differences are bounded in Jarzynski-type
equalities and estimated in annealed importance sampling. That analogy is what
supplies the training objective, and it is the reason the field speaks of
temperature, annealing and entropy at all.

The direction of the forward process is fixed by the second law, which for the
Fokker--Planck equation \eqref{eq:fokker-planck} takes an exact form. Let $p_t$
evolve with a gradient drift and stationary law $p_\infty \propto e^{-V}$, and
consider the free energy $\mathcal F[p_t] = \kl{p_t}{p_\infty}$, which is the
mean energy minus the entropy up to a constant. Writing the dynamics as
$\partial_t p = \nabla \cdot \big(p \nabla \log (p/p_\infty)\big)$ and
integrating by parts,
\begin{equation}
  \frac{\dd \mathcal F}{\dd t}
  = \int \partial_t p \, \log \frac{p}{p_\infty}
  = -\int p \,\Big\| \nabla \log \frac{p}{p_\infty} \Big\|^2 \;\le\; 0 ,
  \label{eq:h-theorem}
\end{equation}
with equality only at $p_t = p_\infty$. This is the H-theorem: the free energy
decreases monotonically and the relative entropy is consumed at a rate given by
a relative Fisher information. The forward channel therefore destroys
information about the data monotonically and irreversibly, and
\eqref{eq:h-theorem} identifies the rate. The reverse direction is not
forbidden, but it is not free either: by Anderson's theorem it requires the
score of the intermediate laws to be supplied from outside. The entire
modelling effort of a diffusion model is the estimation of that one object, and
Chapters~\ref{ch:gaussian} to~\ref{ch:ring} are concerned with what that object
looks like when the data are trajectories.

\section{Denoising diffusion probabilistic models}\label{sec:ddpm}

The discrete-time formulation \citep{sohl2015deep, ho2020denoising}
discretises \eqref{eq:forward-marginal} into a Markov chain
$q(x_s | x_{s-1}) = \Nd(x_s; \sqrt{1 - \beta_s}\, x_{s-1}, \beta_s I)$
with a variance schedule $\beta_1, \dots, \beta_S$, and learns a reverse
chain $p_\theta(x_{s-1} | x_s)$. Training maximises a variational lower
bound on the data log-likelihood; the physicist will recognise the Gibbs
variational principle \eqref{eq:gibbs-variational} in modern dress.

\begin{toolbox}{The evidence lower bound of a diffusion model}
For any latent-variable model with joint $p_\theta(x_{0:S})$ and any
inference distribution $q(x_{1:S} | x_0)$, Jensen's inequality gives
\begin{align*}
  \log p_\theta(x_0)
  &= \log \int p_\theta(x_{0:S})\;\dd x_{1:S}
   = \log\, \E_{q}\!\left[
      \frac{p_\theta(x_{0:S})}{q(x_{1:S}|x_0)} \right] \\
  &\geq\; \E_{q}\!\left[
      \log \frac{p_\theta(x_{0:S})}{q(x_{1:S}|x_0)} \right]
  \;=:\; \mathrm{ELBO}.
\end{align*}
Insert the two Markov factorisations
$p_\theta(x_{0:S}) = p(x_S) \prod_s p_\theta(x_{s-1}|x_s)$ and
$q(x_{1:S}|x_0) = \prod_s q(x_s|x_{s-1})$, then regroup telescopically
using Bayes' rule
$q(x_s|x_{s-1}) = q(x_{s-1}|x_s, x_0)\, q(x_s|x_0) / q(x_{s-1}|x_0)$:
\begin{align*}
  -\mathrm{ELBO}
  \;=\;& \kl{q(x_S|x_0)}{p(x_S)}
  \;-\; \E_q \log p_\theta(x_0 | x_1) \\
  &+ \sum_{s=2}^{S}
    \E_q\, \kl{q(x_{s-1}|x_s, x_0)}{p_\theta(x_{s-1}|x_s)} ,
\end{align*}
one denoising step per summand. The decisive point: because the forward
chain is Gaussian and conditioned on \emph{both} endpoints, each target
$q(x_{s-1} | x_s, x_0)$ is an explicit Gaussian. Matching it with a
Gaussian $p_\theta$ reduces each KL term to a weighted regression of the
posterior mean; concretely, to predicting the noise $\varepsilon$ that
was added \citep{ho2020denoising}. The thermodynamic reading of the ELBO
gap as an entropy-production bound goes back to \citet{sohl2015deep}.
\hfill$\square$
\end{toolbox}

The practical upshot of \citet{ho2020denoising} is the loss
\begin{equation}
  \mathcal{L}(\theta)
  = \E_{s,\, a \sim p_0,\, \varepsilon}
    \big\| \varepsilon - \varepsilon_\theta\big(
      \underbrace{a\,e^{-t_s} + \sqrt{\Delta_{t_s}}\,\varepsilon}_{x_{t_s}},\; s
    \big) \big\|^2 ,
  \label{eq:ddpm-loss}
\end{equation}
a denoising regression, and the empirical discovery that this simplified
objective, with each noise level weighted equally and all the variational
bookkeeping dropped, trains photorealistic image models. Its connection
with the score is not incidental, as the next section shows: predicting
the noise \emph{is} estimating $\nabla \log p_t$ up to scale.

\section{Score matching}\label{sec:score-matching}

Suppose we model the score directly by a function $s_\theta(x)$ and wish
it to match $\nabla \log p(x)$ for data from $p$. The naive objective
$\E_p \| s_\theta - \nabla \log p \|^2$ seems to require knowing the very
score we lack. Two classical results remove the obstruction.

\begin{toolbox}{Implicit score matching (Hyv\"arinen)}
Expand the square and keep only $\theta$-dependent terms:
\begin{equation*}
  \E_p \norm{s_\theta - \nabla \log p}^2
  = \E_p \norm{s_\theta}^2
  - 2\, \E_p \inner{s_\theta}{\nabla \log p}
  + \text{const}.
\end{equation*}
The cross term integrates by parts. Componentwise, using
$p\,\partial_i \log p = \partial_i p$ and assuming
$p\, s_{\theta,i} \to 0$ at infinity:
\begin{equation*}
  \int p\; s_{\theta,i}\; \partial_i \log p \;\dd x
  = \int s_{\theta,i}\; \partial_i p \;\dd x
  = - \int p\; \partial_i s_{\theta,i} \;\dd x .
\end{equation*}
Therefore
\begin{equation*}
  \E_p \norm{s_\theta - \nabla \log p}^2
  = \E_p\Big[ \norm{s_\theta(x)}^2
    + 2\, \nabla \cdot s_\theta(x) \Big] + \text{const} :
\end{equation*}
an objective involving only the model and data samples
\citep{hyvarinen2005estimation}. Elegant, but the divergence term is
expensive for large networks, which is why diffusion models use the
denoising variant instead. \hfill$\square$
\end{toolbox}

\begin{toolbox}{Denoising score matching (Vincent)}
Fix $t$ and consider the noised marginal \eqref{eq:pt-convolution},
written as $p_t(x) = \int p_0(a)\, q_t(x|a)\,\dd a$ with Gaussian kernel
$q_t(x|a) = \Nd(x; a e^{-t}, \Deltat I)$. Then
\begin{equation*}
  \nabla_x \log p_t(x)
  = \frac{\int p_0(a)\, \nabla_x q_t(x|a)\, \dd a}{p_t(x)}
  = \int \underbrace{\frac{p_0(a)\, q_t(x|a)}{p_t(x)}}_{=\,p(a|x)}
    \nabla_x \log q_t(x|a)\, \dd a ,
\end{equation*}
so that
\begin{equation*}
  \nabla_x \log p_t(x)
  \;=\; \E\big[ \nabla_x \log q_t(x|a) \,\big|\, x \big] :
\end{equation*}
the score of the \emph{marginal} is the posterior average of the score of
the \emph{kernel}, and the latter is trivial,
$\nabla_x \log q_t(x|a) = -(x - a e^{-t})/\Deltat$. A standard
$L^2$-projection argument turns this into a regression: for any
$s_\theta$,
\begin{equation*}
  \E_{x \sim p_t} \norm{ s_\theta(x) - \nabla \log p_t(x) }^2
  = \E_{a, x} \norm{ s_\theta(x) - \nabla_x \log q_t(x|a) }^2
  + \text{const},
\end{equation*}
because conditional expectation is exactly the $L^2$ projection onto
functions of $x$ \citep{vincent2011connection}. Substituting the Gaussian
kernel score and $x = a e^{-t} + \sqrt{\Deltat}\varepsilon$ gives
$\nabla_x \log q_t = -\varepsilon / \sqrt{\Deltat}$: \emph{learning the
score of noisy data is learning to predict the noise}, and the DDPM loss
\eqref{eq:ddpm-loss} is denoising score matching in disguise.
\hfill$\square$
\end{toolbox}

The continuous-time synthesis \citep{song2019generative,
song2021scorebased} trains one network $s_\theta(x, t)$ against denoising
score matching at all noise levels and samples with the reverse SDE
\eqref{eq:reverse-sde}; an accessible exposition is \citet{song2021blog},
and \citet{karras2022elucidating} systematise the design space. This is
the resolution, promised in Section~\ref{sec:ebm-def}, of the
energy-based model's oldest problem: \citeauthor{song2019generative}
replace the energy by its gradient, the equilibrium sampler by a
finite-time reverse SDE, and the intractable likelihood by a regression
whose targets are exact.

\section{Flows: the deterministic side of the same coin}\label{sec:flows}

Anderson's theorem has a deterministic twin. As noted in
Section~\ref{sec:reversal}, the marginals $p_t$ of the forward SDE are
also transported by the noise-free \emph{probability-flow ODE} with
velocity $f - \tfrac{g^2}{2}\nabla \log p_t$ \citep{song2021scorebased}:
the same score, consumed by an ODE instead of an SDE, turns generation
into a smooth, invertible map from Gaussian noise to data. This places
diffusion models inside an older family. \emph{Normalizing flows}
\citep{rezende2015variational, papamakarios2021normalizing} construct an
invertible map $T_\theta$ and evaluate likelihoods through the
change-of-variables formula; their continuous-time version parametrises
the velocity field of an ODE directly \citep{chen2018neural}. What
diffusion brought to this family is a \emph{simulation-free training
signal}: \emph{flow matching} \citep{lipman2023flow} and stochastic
interpolants \citep{albergo2023building} show that one can regress the
velocity field of any prescribed interpolation between noise and data
(the OU path \eqref{eq:forward-marginal} being one choice among many)
with the same conditional trick as denoising score matching, never
integrating the ODE during training; a systematic introduction is
\citet{holderrieth2025flow}. The unifying statement is worth recording:
DDPMs, score SDEs, and flow-matching models all learn a vector field
whose ground truth is a \emph{conditional expectation under the noising
channel}; they differ in the path of marginals and in whether sampling is
stochastic or deterministic. Every exact result of this thesis therefore
speaks to the whole family at once: the object we compute in closed form,
the conditional expectation behind the field, is the training target of
all of them.

\section{The score as a posterior expectation}\label{sec:tweedie}

The middle line of Vincent's derivation deserves its own section,
because every result of this thesis passes through it. Rearranging
$\nabla \log p_t(x) = \E[-(x - a e^{-t})/\Deltat \mid x]$:
\begin{equation}
  \boxed{\;
  \nabla_x \log p_t(x)
  \;=\; \frac{e^{-t}\,\E[\,a \,|\, x\,] - x}{\Deltat}
  \quad\Longleftrightarrow\quad
  \E[\,a \,|\, x\,]
  = e^{t} \Big( x + \Deltat\, \nabla_x \log p_t(x) \Big).
  \;}
  \label{eq:tweedie}
\end{equation}
We refer to \eqref{eq:tweedie} throughout as the \emph{score--posterior
identity}; in the empirical-Bayes literature it is known as Tweedie's
formula \citep{robbins1956empirical, miyasawa1961empirical,
efron2011tweedie}. The score of the noisy marginal and the Bayes-optimal
denoiser $\E[a|x]$ determine each other exactly, with no assumption
whatsoever on $p_0$. Three readings will recur:
\begin{enumerate}
  \item \emph{Geometric}: the score points from $x$ towards (the
  contracted image of) the posterior mean of the clean data; denoising
  direction and score direction coincide.
  \item \emph{Statistical}: estimating the score is solving a Bayesian
  inference problem, namely computing a posterior expectation. The score
  is not \emph{like} an inference target; it \emph{is} one.
  \item \emph{Algorithmic}: for sequence data with a Markov prior, that
  inference problem lives on a tree-shaped graph, so the score is
  computable by message passing, the programme made precise in
  Chapter~\ref{ch:graphs} and executed in Chapter~\ref{ch:gaussian}.
\end{enumerate}

For \emph{vector-valued} data $a \in \R^K$ noised coordinate-wise,
\eqref{eq:tweedie} holds componentwise with the \emph{full} posterior:
\begin{equation}
  \big( \nabla_x \log p_t(x) \big)_k
  \;=\; \frac{e^{-t}\, \E[\,a_k \,|\, x_0, \dots, x_{K-1}\,] - x_k}{\Deltat}.
  \label{eq:tweedie-joint}
\end{equation}
The conditioning is on \emph{all} coordinates: even though the noise is
independent across frames, the posterior mean of frame $k$, and hence the
$k$-th score component, depends on every observation, because the prior
couples the frames. Equation \eqref{eq:tweedie-joint} is the precise
reason the \emph{joint} score differs from the per-frame marginal score,
a distinction that is easy to blur and whose importance for structured
data was impressed on this project in supervision
(Chapter~\ref{ch:development}); Chapter~\ref{ch:gaussian} quantifies
exactly how much coupling the conditioning induces, as a function of $t$.

\section{The statistical physics of generative diffusion}
\label{sec:statphys-diffusion}

A line of work initiated by Biroli, M\'ezard, and collaborators analyses
diffusion models with the tools of Chapter~\ref{ch:statmech}, in the
regime where the score is \emph{exact}: either computed from a solvable
$p_0$ or from the empirical measure of $n$ training points. Because this
is the literature our research extends, we review it in some detail.

\paragraph{High dimensions and the cost of asymmetry.}
\citet{biroli2023generative} analyse generative diffusion in the
high-dimensional limit for two controlled data models, a Gaussian model
and the Curie--Weiss ferromagnet. In the latter, the reverse dynamics
undergoes the symmetry-breaking mechanism of Section~\ref{sec:ising}: at
a critical diffusion time the trajectory commits to one of the two
low-temperature states. Their sharpest quantitative point concerns data
requirements: to reconstruct the \emph{relative weights} of two
asymmetric modes, the number of training samples must be much larger than
the data dimension, and the paper characterises the scaling laws in
dimension and sample size for efficient generation.

\paragraph{The three regimes of the backward dynamics.}
\citet{biroli2024dynamical} study the backward process when the score has
been trained optimally on $n$ samples in dimension $d$, and identify
three dynamical regimes separated by two transitions. Starting from pure
noise, the trajectory first behaves as a free (Brownian-like) process;
at the \emph{speciation} time it commits to the gross structure of the
data (a class, a mode), through a mechanism analogous to spontaneous
symmetry breaking; at the later \emph{collapse} time the trajectory is
attracted to an individual memorised training point, through a mechanism
analogous to condensation in a glass. The speciation time is computable
from a spectral analysis of the data correlation matrix, and the collapse
time from an excess-entropy estimate; the dependence of the collapse time
on $n$ and $d$ quantifies the curse of dimensionality for diffusion
models. Related transition phenomena along the generation trajectory were
exhibited by \citet{raya2023spontaneous}.

\paragraph{Speciation with general class structure.}
\citet{achilli2026speciation} develop the general theory of the
speciation transition: where earlier analyses required classes separated
by their first moments (Gaussian mixtures with distinct means), they
characterise speciation times for arbitrary class structures through a
free-entropy difference criterion, recovering the known Gaussian-mixture
results and extending them to classes that differ only through
higher-order or collective features, with successive speciation events
for multiple classes. A hierarchical counterpart is
\citet{sclocchi2025phase}: in a hierarchical generative model of data,
the backward process started from time $t$ exhibits a phase transition at
a threshold time at which the probability of preserving the high-level
class drops abruptly, while low-level features evolve smoothly; beyond
the transition a regenerated sample can retain low-level elements of the
original image while its class has changed. The statistical-physics
treatment of generative diffusion in solvable models is developed at
book length in the doctoral thesis of \citet{achilli2026thesis}.

\paragraph{Memorisation and generalisation.} When the score is learned
from finitely many samples rather than computed exactly, the question
becomes when generation generalises and when it memorises.
\citet{bonnaire2025memorize} identify two timescales in the training
dynamics: an early time $\tau_{\mathrm{gen}}$ at which the model begins
to generate high-quality samples, and a later time $\tau_{\mathrm{mem}}$
beyond which memorisation of training data emerges;
$\tau_{\mathrm{mem}}$ grows linearly with the training-set size $n$ while
$\tau_{\mathrm{gen}}$ stays constant, opening a widening window in which
overparameterised models generalise, a form of implicit dynamical
regularisation. \citet{garnierbrun2026biased} refine the picture by
identifying a phase of \emph{biased} generalisation in between: the test
loss still decreases while the model increasingly favours samples
anomalously close to training points, so the standard early-stopping
criterion does not certify genuine novelty. \citet{ghio2024sampling} map
generative samplers (flows, diffusions, autoregressive networks) onto
statistical-physics sampling problems and delimit where each is
efficient, importing the algorithmic-hardness picture of spin glasses
\citep{zdeborova2016statistical}.

\paragraph{The exactly solvable programme, and what it holds fixed.} The
common method of these works is to choose $p_0$ rich enough to be
interesting and simple enough that $\nabla \log p_t$ is available in
closed form, then study the exact generation dynamics. What none of them
vary is the \emph{internal structure of a single data point}: their
coordinates are exchangeable or i.i.d.\ conditional on a mode. There is
no ordering and no dynamics \emph{inside} the sample. That axis is
precisely where this thesis works.

\section{Learning as a cascade of phase transitions}\label{sec:cascade}

The works above put phase transitions on the \emph{sampling} side of
generative modelling. A complementary result, \citet{bachtis2024cascade},
which this section examines in detail, puts them on the \emph{training}
side, and completes the argument of Section~\ref{sec:statmech-learning}
that learning itself is a sequence of symmetry breakings. The object of
study is the restricted Boltzmann machine of
Section~\ref{sec:boltzmann-machines}, the minimal trainable EBM, with
energy $E(v, h) = -v^\T W h - b^\T v - c^\T h$ and the moment-matching
gradient \eqref{eq:bm-learning}. The central idea is to track training
through the singular value decomposition of the weight matrix,
\begin{equation}
  W \;=\; \sum_{\gamma=1}^{r} w_\gamma\, u_\gamma \bar u_\gamma^\T ,
  \label{eq:rbm-svd}
\end{equation}
whose modes $(u_\gamma, \bar u_\gamma)$ act as candidate order
parameters: the magnetisation of the visible layer along mode $\gamma$,
$m_\gamma = u_\gamma \cdot v / \sqrt{N_v}$, measures whether the learned
Gibbs measure has developed structure in that direction. At
initialisation the weights are small and the model is a paramagnet: no
direction in configuration space is preferred. The claim of
\citet{bachtis2024cascade} is that training is an \emph{effective
cooling}: as the singular values grow, the model crosses a sequence of
genuine second-order phase transitions, one for each feature it learns,
each signalled by a diverging susceptibility and falling in the
mean-field (Curie--Weiss) universality class. The mechanism is visible in
a two-line calculation.

\begin{toolbox}{The retrieval transition of a one-mode RBM}
Following \citet{bachtis2024cascade}, take a single Gaussian hidden unit
of variance $1/N_v$ coupled to $N_v$ binary visibles through a weight
vector $w$: the energy is
$E(v, h) = -h\, w \cdot v + \tfrac{N_v}{2} h^2$. Integrating out the
Gaussian $h$ gives the visible marginal
\begin{equation*}
  p(v) \;\propto\; \exp\!\Big[ \tfrac{1}{2 N_v}\, (w \cdot v)^2 \Big] :
\end{equation*}
a Mattis/Hopfield measure (Section~\ref{sec:hopfield}) whose single
pattern is $w$ itself, made extensive by the $1/N_v$ variance scaling.
Write $w = \bar w\, \xi$ with $\xi \in \{-1,+1\}^{N_v}$ and let
$m = \tfrac{1}{N_v} \sum_i \xi_i v_i$ be the overlap; then
$p(v) \propto \exp[\tfrac{N_v \bar w^2}{2} m^2]$. Counting configurations
at fixed overlap ($\asymp \exp[N_v H(\tfrac{1+m}{2})]$ with $H$ the
binary entropy) gives the free-entropy density
\begin{equation*}
  \phi(m) \;=\; \frac{\bar w^2}{2}\, m^2 + H\!\Big(\frac{1+m}{2}\Big)
  \;=\; \log 2 + \frac{\bar w^2 - 1}{2}\, m^2 - \frac{m^4}{12} + O(m^6).
\end{equation*}
This is a Landau expansion: the quadratic coefficient changes sign at
$\bar w^2 = \norm{w}^2 / N_v = 1$. Below the threshold the maximiser is
$m = 0$ (paramagnet); above it two symmetric maximisers $\pm m^*$ appear
continuously (retrieval). The transition is second order with mean-field
exponents, and the susceptibility along $w$ diverges as
$\norm{w}^2/N_v \uparrow 1$. Growing a singular value during training is
therefore literally a quench through a Curie--Weiss critical point.
\hfill$\square$
\end{toolbox}

\paragraph{Why a cascade.} With data generated by a Hopfield teacher
whose two patterns are correlated combinations
$\xi^{1,2} = \eta_1 \pm \eta_2$ of orthogonal directions, the data
covariance is rank two with unequal strengths, and the early-time
gradient flow of the student's SVD amplitudes along $\eta_1$ and $\eta_2$
is exponential with \emph{two different rates}. The stronger direction
(the centre of mass of the clusters) reaches its critical magnitude
first: the learned measure splits into two lumps along $\pm\eta_1$. The
weaker direction keeps growing at its own rate and crosses criticality
later, splitting each lump again: four modes, the full teacher, retrieved
in hierarchical order. Each crossing is a distinct second-order
transition with its own critical time, its own diverging susceptibility,
and hysteresis below it. Learning proceeds coarse-to-fine not by accident
of optimisation but because the critical times are ordered by the
spectrum of the data.

\paragraph{Evidence on real data.} On MNIST, CelebA, and human-genome
data, \citet{bachtis2024cascade} track the susceptibilities of individual
SVD modes during training of binary--binary RBMs and observe the same
signatures: successive susceptibility peaks that sharpen with system
size, finite-size-scaling collapses consistent with mean-field exponents
$(\gamma, \nu) = (1, \tfrac12)$ at upper critical dimension four, and
hysteresis, the operational fingerprints of genuine transitions rather
than smooth crossovers. This confirms and sharpens the earlier picture of
RBM learning as progressive alignment with the data's principal
components \citep{decelle2017spectral}, of feature emergence as
condensation \citep{tubiana2017emergence}, and of sequential mode
learning in linear networks \citep{saxe2014exact}; a broad review of the
EBM--physics interface is \citet{carbone2025hitchhiker}.

\paragraph{What this means for the thesis.} Phase transitions now appear
on both faces of generative diffusion: in \emph{sampling}, as speciation
events along the reverse trajectory \citep{biroli2024dynamical,
achilli2026speciation}, and in \emph{training}, as the cascade just
described. Both are hierarchies of symmetry breakings ordered by the
data's spectrum. Our research adds the axis that both literatures hold
fixed: the structure of the inference problem \emph{inside a single data
point}, and how the diffusion time $t$ deforms it. The chain we solve is
the minimal model in which that question is non-trivial and answerable
exactly.

\section{The gap, restated}\label{sec:gap}\label{sec:related}

Real sequential data (video, audio, trajectories) have exactly the
structure the solvable programme leaves out: an internal (frame) time
with a Markovian dependence along it. For such data the score
\eqref{eq:tweedie-joint} is a smoothing posterior on a chain, an object
with seventy years of inference theory behind it. Yet the questions a
diffusion modeller asks about it are not answered by the classical
literature, which fixes the noise level, nor by the diffusion literature,
which ignores internal structure: how does the coupling structure of
$\nabla \log p_t$ deform with $t$? Can it be computed locally? What
survives a Gaussian parametrisation of the messages when the data are not
Gaussian? What limited guidance does that give for score architectures?
Answering these questions, in the fully solvable setting of the AR(1)
chain and its non-Gaussian perturbations, is the contribution of
Chapters~\ref{ch:gaussian} and~\ref{ch:laplace}, with the inference tools
prepared in Chapter~\ref{ch:graphs} and the path that led to the precise
formulation recorded in Chapter~\ref{ch:development}.
